\chapter{Conclusions}

\label{chap:conclusions}

This thesis investigated the modeling, design, and control of an Aerial Cable-Towed System (ACTS) intended for collaborative payload manipulation with focus on understanding how the system's architecture influences its ability to generate the required payload wrench while maintaining stable and accurate motion.

A general kinematostatic model of the ACTS was first developed by extending concepts from CDPRs to hybrid ACTSs. After that, a centralized control architecture was implemented in simulation, which was able to compute the desired payload wrench and command the corresponding UAV commands. A simulation framework based on MuJoCo was developed to evaluate the closed-loop behavior of the selected architectures under identical operating conditions.

Once the controller was proven to work, a screening procedure was proposed to efficiently explore the large design space of possible cable routings, analyzing feasibility constraints, such as WCC, IFC and WFC, and geometric performance indices, such as conditioning index and manipulability.

The simulation results demonstrate that the proposed pre-screening successfully removes many unsuitable cable routings while retaining architectures capable of achieving satisfactory dynamic performance. At the same time, they highlight that geometric performance indices alone are insufficient to predict the behavior of the complete closed-loop system. Architectures exhibiting excellent conditioning and manipulability do not necessarily achieve the highest pose reachability or the lowest tracking errors once payload dynamics, cable tension redistribution, and controller interaction are taken into account. Dynamic simulation therefore represents an essential second stage of the design process, complementing the theoretical analysis rather than replacing it. Perfect example is \texttt{rectangle-opposite-2x3-triangle} which tops every static metric, yet has the worst pose reachability; differently, \texttt{pentagon-rectangle-same-triangle} achieve the best pose reachability and the reference configuration (\texttt{acts\_stewart}), last on every static index, has the lowest cable rubbing rate.

Secondly, a more spread-out attachment point on the ground consistently provides higher performances with better indices and fewer unreached poses. However, this comes with a larger installation area on the deployment site and a need to find a compromise between the two.

Lastly, the simulations carried on demonstrate how the choice of the drone attachments to the payload is crucial to successfully reach the desired payload pose. In fact, while linear disposition demonstrates in theory large potential, only the triangular disposition provides independent moment arms in two dimensions, substantially improving the task completion rate.

The proposed approach reduces the complexity of the architectural design process while providing quantitative criteria for selecting cable configurations suitable for collaborative aerial manipulation. However, several aspects remain open for future investigation. 

Firstly, using points belonging to the desired payload trajectory rather than random ones may further improve the architecture selection of the system.

Secondly, while this thesis focuses on the analysis of six ground cables and three drones, future work may extend the analysis to different combinations of ground and aerial actuators, answering the question of how their number affects may extend the analysis to different combinations of ground and aerial actuators, answering the question of how their number affects the system performance. Additionally, the connection strategy between drones and payload may also be analyzed deeply, investigating whether connecting each UAV through one or two cables provides advantages in terms of maneuverability, feasible workspace, and available wrench set. These considerations, together with considering aerial and ground cables connected to arbitrary payload faces, assess their influence on workspace and practical usability.

Another natural progression is the inclusion of physical interaction with vertical surfaces to assess the influence of contact forces. 

More importantly, the carried-out methodology should be validated experimentally on a real ACTS platform to assess the influence of sensing uncertainties and unmodeled aerodynamic effects. The current simulations neglect cable elasticity, rotor wash, and environmental disturbances, whose inclusion would improve the realism of the model.